\section{Marco teórico}
\label{sec:marco}

\subsection{Ecuaciones}

Las ecuaciones numeradas se referencian con \verb|\cref|. Por ejemplo, la entropía de un
conjunto $S$ con $c$ clases (\cref{eq:entropia}) y la ganancia de información de un
atributo $A$ (\cref{eq:ganancia}):
\begin{equation}
  H(S) = -\sum_{i=1}^{c} p_i \log_2 p_i
  \label{eq:entropia}
\end{equation}
\begin{equation}
  \mathrm{IG}(S, A) = H(S) - \sum_{v \in \mathrm{Valores}(A)} \frac{\abs{S_v}}{\abs{S}}\, H(S_v)
  \label{eq:ganancia}
\end{equation}

Los comandos definidos en \texttt{informe.sty} simplifican la notación habitual:
$\vect{x} \in \R^{d}$, $\mat{X} \in \R^{n \times d}$, $\norm{\vect{x} - \vect{\mu}}_2$,
$\E[X]$, $\Var(X)$ y
\begin{equation*}
  \hat{\vect{\beta}} = \argmin_{\vect{\beta}} \norm{\vect{y} - \mat{X}\vect{\beta}}_2^2 .
\end{equation*}

\subsection{Recuadros}

\begin{nota}
  Útil para definiciones, supuestos o aclaraciones que conviene destacar del texto.
\end{nota}

\begin{importante}[Restricción de integridad]
  En MySQL 8.4 las restricciones \texttt{CHECK} se validan efectivamente; en versiones
  anteriores a la 8.0.16 se ignoraban.
\end{importante}
